\documentclass[11pt]{article}
\usepackage{amsmath,amssymb,amsthm}
\usepackage[margin=1in]{geometry}
\usepackage{mathtools}
\usepackage[hidelinks]{hyperref}

\newtheorem{theorem}{Theorem}
\newtheorem{lemma}[theorem]{Lemma}
\newtheorem{definition}[theorem]{Definition}
\newtheorem{corollary}[theorem]{Corollary}
\newtheorem{remark}[theorem]{Remark}

%% -- Notation --
\newcommand{\pair}[2]{\langle #1, #2 \rangle}
\newcommand{\trueT}{\mathbf{t}}
\newcommand{\falseT}{\mathbf{f}}
\newcommand{\SchemaF}{\textbf{Schema~F}}
\newcommand{\nat}[1]{\overline{#1}}
\newcommand{\code}[1]{\ulcorner #1 \urcorner}
\newcommand{\reify}[1]{\widehat{#1}}
\newcommand{\thFun}{\mathrm{thFun}}
\newcommand{\thFunTFor}{\mathrm{thFunTFor}}
\newcommand{\cstf}{\mathrm{cstf}}
\newcommand{\TreeEq}{\mathrm{TreeEq}}
\newcommand{\IfLf}{\mathrm{IfLf}}

\title{G\"odel's Second Incompleteness Theorem\\
  in a Guard/Rose Equational System:\\
  Formalization in Agda}
\author{}
\date{}

\begin{document}
\maketitle

%% ===============================================================
\section{Overview}

We formalize G\"odel's first and second incompleteness theorems
in a purely equational system combining two ideas:
\begin{itemize}
\item \emph{Guard}'s syntax~\cite{guard1963}: binary trees with
  typed unary and binary functors, serving as both the data domain
  and the term language;
\item \emph{Rose}'s purely equational proof style~\cite{rose1961}:
  the only judgement is an equation $t = u$ between terms, with no
  propositional connectives, no sequents, and no logical rules ---
  only computation axioms, congruence, substitution, and a
  tree-induction principle (\SchemaF).
\end{itemize}
The formalization is in Agda with
\texttt{-{}-safe -{}-without-K -{}-exact-split}.  No postulates.

\begin{theorem}[G\"odel~I, constructive form --- \texttt{godelIToBot} in \texttt{ProofEAny.agda}]
\label{thm:godel1c}
For any hypothesis equation~$H$:
\[
  H \vdash G \;\to\; H \vdash \trueT = \falseT
\]
where $G$ is the G\"odel sentence.
\end{theorem}

\begin{theorem}[G\"odel~II, constructive form --- \texttt{conToBot} in \texttt{GodelIIFull.agda}]
\label{thm:godel2c}
For any hypothesis equation~$H$:
\[
  H \vdash \mathrm{Con} \;\to\;
  H \vdash \trueT = \falseT.
\]
\end{theorem}

\noindent
Both results have standard corollaries:
if~$H$ is consistent, then neither~$G$ nor $\mathrm{Con}$
is derivable from~$H$.

%% ===============================================================
\section{The formal system}

\subsection*{Trees}

The ground data are finite binary trees:
\[
  a, b \;::=\; O \;\mid\; \pair{a}{b}
\]
where $O$ is the leaf and $\pair{a}{b}$ is the node.
We define:
\[
  \trueT \;=\; O, \qquad \falseT \;=\; \pair{O}{O}.
\]
Here $\trueT$ represents equality (or ``yes'') and $\falseT$
represents inequality (or ``no'').
We write $\nat{n}$ for the unary encoding of $n \in \mathbb{N}$:
$\nat{0} = O$, $\nat{n+1} = \pair{O}{\nat{n}}$.

\subsection*{Terms and functors}

Terms are built from trees, variables, and functor applications.
We use $t, u, v$ for terms, $a, b$ for trees.
The system has countably many variables $x_0, x_1, x_2, \ldots$\,.
There are unary functors ($f, g : \mathrm{Fun}_1$) and binary
functors ($h : \mathrm{Fun}_2$).

\[
\begin{array}{rcl@{\qquad}l}
  f, g &::=& \mathrm{I} & \text{identity} \\
       &\mid& \pi_1 & \text{first projection} \\
       &\mid& \pi_2 & \text{second projection} \\
       &\mid& f \circ g & \text{composition} \\
       &\mid& h[f, g] & \text{binary composition: } t \mapsto h(f(t), g(t)) \\
       &\mid& \mathrm{R}(z, s) & \text{tree recursion with base } z \text{ and step } s \\
       &\mid& \mathrm{K}_c & \text{constant } c \\[6pt]
  h &::=& \mathrm{P} & \text{pairing: } (t, u) \mapsto \pair{t}{u} \\
    &\mid& \mathrm{C} & \text{left projection: } (t, u) \mapsto t \\
    &\mid& \mathrm{Lf}(f) & \text{lift: } (t, u) \mapsto f(t) \\
    &\mid& \mathrm{Ps}(f, h) & \text{post-compose: } (t, u) \mapsto f(h(t, u)) \\
    &\mid& \mathrm{Fn}(h_1, h_2, h) & \text{fan: } (t, u) \mapsto h(h_1(t,u),\, h_2(t,u)) \\
    &\mid& \mathrm{If} & \text{conditional} \\
    &\mid& \TreeEq & \text{tree equality test}
\end{array}
\]

Terms:
\[
  t, u \;::=\; O \;\mid\; x_i \;\mid\; f(t) \;\mid\; h(t, u)
\]

An \emph{equation} is a pair $t = u$ of terms.

\subsection*{Embedding trees into terms}

Every tree $a$ determines a closed term $\hat{a}$:
\[
  \hat{O} = O, \qquad \widehat{\pair{a}{b}} = \pair{\hat a}{\hat b}.
\]
(In Agda: \texttt{reify}.)
Conversely, every closed term reduces to $\hat{a}$ for a unique~$a$.

\subsection*{Derivation rules}

A derivation is built from the rules below.
Every rule may use a hypothesis~$H$;
we write $H \vdash t = u$.

\medskip\noindent
\textbf{Computation axioms.}
\begin{alignat*}{2}
  &\text{Identity:}          &\quad \mathrm{I}(t) &= t \\
  &\text{First projection:}  &\quad \pi_1(\pair{t}{u}) &= t \\
  &\text{Second projection:} &\quad \pi_2(\pair{t}{u}) &= u \\
  &\text{Constant:}          &\quad \mathrm{C}(t, u) &= t \\
  &\text{Composition:}       &\quad (f \circ g)(t) &= f(g(t)) \\
  &\text{Binary composition:}&\quad h[f,g](t) &= h(f(t),\, g(t)) \\
  &\text{Lift:}              &\quad \mathrm{Lf}(f)(t, u) &= f(t) \\
  &\text{Post-compose:}      &\quad \mathrm{Ps}(f, h)(t, u) &= f(h(t,u)) \\
  &\text{Fan:}               &\quad \mathrm{Fn}(h_1,h_2,h)(t,u)
                                      &= h(h_1(t,u),\, h_2(t,u)) \\
  &\text{Constant functor:}  &\quad \mathrm{K}_c(t) &= c
\end{alignat*}

\noindent
\textbf{Tree recursion.}
\begin{alignat*}{2}
  &\text{Base:} &\quad \mathrm{R}(z, s)(O) &= z \\
  &\text{Step:} &\quad \mathrm{R}(z, s)(\pair{t}{u})
    &= s\bigl(\pair{t}{u},\;
         \pair{\mathrm{R}(z,s)(t)}{\mathrm{R}(z,s)(u)}\bigr)
\end{alignat*}

\noindent
\textbf{Conditional.}
\begin{alignat*}{2}
  &\text{True branch:}     &\quad \mathrm{If}(\trueT,\, \pair{t}{u}) &= t \\
  &\text{False branch:}     &\quad \mathrm{If}(\pair{t_1}{t_2},\, \pair{u_1}{u_2}) &= u_2
\end{alignat*}
More generally, $\mathrm{If}$ branches on whether its first
argument is~$O$ (leaf) or a pair (node).

\medskip\noindent
\textbf{Tree equality.}
\begin{alignat*}{2}
  &\TreeEq(O, O) &\;=\; \trueT \\
  &\TreeEq(O, \pair{t}{u}) &\;=\; \falseT \\
  &\TreeEq(\pair{t}{u}, O) &\;=\; \falseT \\
  &\TreeEq(\pair{t_1}{t_2}, \pair{u_1}{u_2})
    &\;=\; \mathrm{If}\bigl(\TreeEq(t_1, u_1),\;
         \pair{\TreeEq(t_2, u_2)}{\falseT}\bigr)
\end{alignat*}
$\TreeEq$ returns $\trueT$ for equal trees and $\falseT$ for
unequal.  The last clause says: if the left children differ, return
$\falseT$ immediately; otherwise, recurse on the right children.

In particular $\TreeEq(t, t) = \trueT$ for all~$t$,
proved by \SchemaF{} (\texttt{treeEqSelf} in \texttt{GodelII.agda}).

\medskip\noindent
\textbf{Structural rules.}
\begin{gather*}
  \frac{\phantom{t=t}}{t = t}\;\text{(refl)}
  \qquad\qquad
  \frac{t = u}{u = t}\;\text{(sym)}
  \qquad\qquad
  \frac{t = u \qquad u = v}{t = v}\;\text{(trans)}
  \\[8pt]
  \frac{t = u}{f(t) = f(u)}\;\text{(cong$_1$)}
  \qquad\qquad
  \frac{t = u}{h(t, v) = h(u, v)}\;\text{(cong$_L$)}
  \qquad\qquad
  \frac{t = u}{h(v, t) = h(v, u)}\;\text{(cong$_R$)}
\end{gather*}

\noindent
\textbf{Substitution.}
\[
  \frac{H \vdash t = u}
       {H \vdash t[c/x_i] = u[c/x_i]}
  \;\text{(inst)}
\]

\noindent
\textbf{Hypothesis.}
\[
  \frac{\phantom{H}}{H \vdash H}
  \;\text{(hyp)}
\]

\noindent
\textbf{\SchemaF} (tree induction).
If $f$ and $g$ are unary functors, $z$ a term, $s$ a binary functor,
and we have derivations of:
\begin{align*}
  f(O) &= z,  &  f(\pair{x_0}{x_1}) &= s(\pair{x_0}{x_1},\; \pair{f(x_0)}{f(x_1)}), \\
  g(O) &= z,  &  g(\pair{x_0}{x_1}) &= s(\pair{x_0}{x_1},\; \pair{g(x_0)}{g(x_1)}),
\end{align*}
then $f(x_0) = g(x_0)$.

\medskip
\SchemaF{} expresses the uniqueness of tree recursion:
if $f$ and $g$ satisfy the same recurrence (same base~$z$, same step~$s$),
they agree on all inputs.  Restricted to unary trees $\nat{n}$,
this gives ordinary mathematical induction.

\subsection*{Consistency}

The system is \emph{consistent} (relative to hypothesis~$H$)
if there is no derivation of $\trueT = \falseT$ from~$H$.

%% ===============================================================
\section{Key components}

\subsection{Proof encoding ($\thFunTFor$, Theorem~14)}

$\thFunTFor$ is a unary functor
(defined as $\mathrm{R}(O, \mathrm{thFunStep})$)
that maps proof-encoding trees to the code of the equation they prove.
The step function $\mathrm{thFunStep}$ is a 26-branch dispatch
handling all derivation rules.

Every term~$t$ is encoded as a tree $\code{t}$, and every
equation $t = u$ as $\code{t = u} = \pair{\code{t}}{\code{u}}$.
Every functor is similarly encoded.

\emph{Theorem~14} ($\mathrm{thm14E}$): if $H \vdash t = u$,
then there exist trees $p_a, p_b$ such that:
\begin{enumerate}
\item $\thFun(\pair{p_a}{p_b}) = \code{t = u}$ \quad (meta-level correctness),
\item $\vdash \thFunTFor(\pair{\reify{p_a}}{\reify{p_b}}) = \reify{\code{t = u}}$
  \quad (object-level correctness, polymorphic in hypothesis).
\end{enumerate}
The proof is by induction on the derivation, with 26~cases
(1402 lines in \texttt{Thm14E.agda}).

Component~(2) is \emph{polymorphic}: the derivation uses only
computation axioms and structural rules, never the hypothesis
rule.  This is crucial for G\"odel~II.

\subsection{Coded substitution ($\mathrm{substTFor}$)}

$\mathrm{substTFor}$ is a unary functor implementing coded
substitution: given a coded term~$\code{t}$,
it replaces every occurrence of variable~$x_n$ by a
replacement term~$r$.

These two parameters ($r$ and~$n$) are represented by free variables
in $\mathrm{substTFor}$.  To avoid collision with the variables
$x_0, x_1$ used in \SchemaF{} and in the template equation,
the formalization places them at high indices:
\begin{itemize}
\item $x_{11}$ carries the replacement~$r$,
\item $x_{12}$ carries the target index~$\nat{n}$.
\end{itemize}

Closing these variables gives the fully applied version:
\[
  \mathrm{closedSubstTFor}(r, \nat{n}) \;=\;
  \mathrm{substTFor}[r/x_{11},\; \nat{n}/x_{12}].
\]

Correctness: for any term~$t$,
$\mathrm{closedSubstTFor}(r, \nat{n})$ applied to
$\reify{\code{t}}$
equals $\reify{\code{t[r/x_n]}}$.

\subsection{$\TreeEq(t, t) = \trueT$ via \SchemaF}

$\TreeEq(t, t) = \trueT$ for all~$t$ is proved using \SchemaF.
Define $f = \TreeEq[\mathrm{I}, \mathrm{I}]$
(so $f(t) = \TreeEq(t,t)$) and $g = \mathrm{K}_\trueT$.
Both satisfy the same recursion:
\begin{itemize}
\item $f$: by the computation axioms for $\TreeEq$ and the
  recursive structure of $\TreeEq$ on $\pair{t_1}{t_2}$;
\item $g$: by $\mathrm{K}_\trueT$; the step reduces to
  $\mathrm{If}(\trueT, \pair{\trueT}{\falseT}) = \trueT$ since $g$
  returns $\trueT$ on all inputs.
\end{itemize}
\SchemaF{} gives $f(t) = g(t)$, i.e., $\TreeEq(t,t) = \trueT$.

%% ===============================================================
\section{The G\"odel sentence and diagonal construction}

\subsection*{Template}

The G\"odel sentence is built from a \emph{template equation}
with several free variables:
\begin{itemize}
\item $x_0$: will receive a proof encoding (the ``proof slot'');
\item $x_1$: will receive a coded term (the ``diagonal slot'');
\item $x_{11}, x_{12}$: the replacement and target parameters of
  $\mathrm{substTFor}$ (closed later).
\end{itemize}

The template is:
\[
  \mathrm{templateEqn} \;=\;
  \bigl(\TreeEq(\thFunTFor(x_0),\;
    \mathrm{substTFor}(x_1)) = \falseT\bigr).
\]
It says: ``the equation encoded by the proof in~$x_0$ is not equal
to the result of substituting via~$x_1$.''
Leaving $x_{11}, x_{12}$ free at this stage keeps the tree code
of the template small.

\subsection*{Diagonalization}

Let $N = \code{\mathrm{templateEqn}}$ be the tree
encoding of the template.  The G\"odel sentence is obtained by
substituting $\hat{N}$ (the numeral for~$N$) into the diagonal slot:
\[
  G \;=\; \mathrm{templateEqn}[\hat{N}/x_1].
\]
After this substitution, $G$ says (informally):
``for any proof tree in~$x_0$, the equation it encodes is not the
result of substituting~$N$ into the template'' ---
i.e., ``no proof tree proves~$G$.''

\subsection*{Abstract definitions (cached)}

The following are computed once in \texttt{SubstTForPrecomp.agda}
and cached:
\begin{itemize}
\item $N$: tree code of the template
\item $\hat N$: the numeral $\reify{N}$
\item $G = \mathrm{templateEqn}[\hat{N}/x_1]$
\item $\code{G}$: tree encoding of~$G$
\item $\cstf = \mathrm{closedSubstTFor}(\reify{\hat N},\;\nat{1})$:
  the substitution functor closed with replacement = code of the
  template numeral and target = variable~$x_1$
\end{itemize}

\subsection*{Heavy computation (one-time)}

\texttt{SubstTForPrecomp.agda} contains one step requiring
${\sim}108$ seconds (computed once, cached in \texttt{.agdai}):
proving that after instantiating $x_{11}$ and~$x_{12}$, the open
$\mathrm{substTFor}$ becomes~$\cstf$.
This requires Agda to traverse $\reify{\hat{N}}$
(${\sim}400$K term nodes) with substitution acting as identity
on the closed subterm.

\begin{remark}
The mathematical argument does not require this traversal.
A general ``substitution is identity on closed terms'' lemma would
eliminate it.
\end{remark}

%% ===============================================================
\section{Proof of G\"odel I}

\begin{theorem}[G\"odel~I, constructive form]
\label{thm:godel1cf}
For any hypothesis equation~$H$:
$H \vdash G \;\to\; H \vdash \trueT = \falseT$.
\end{theorem}

\begin{corollary}[G\"odel~I, standard form --- \texttt{godelI} in \texttt{ProofEAny.agda}]
\label{thm:godel1s}
$\mathit{Consistent}(H) \;\to\; H \vdash G \;\to\; \bot$.
\end{corollary}

\noindent
Both forms remove the $\mathrm{ProofE}(H)$ assumption
from the original proof.  The key enabler is
$\mathrm{mkProofEAny}$, which constructs a proof encoding
for any equation using the trans-from-two-refls technique.

\medskip\noindent
The proof in \texttt{GodelII.agda} (0.2 seconds, 168 lines)
follows Rose's Theorem~18~\cite{rose1961}.
We assume $H \vdash G$ and derive $H \vdash \trueT = \falseT$:

\begin{enumerate}
\item \textbf{Assume} $H \vdash G$.

\item \textbf{Theorem~14} gives a proof encoding $e$ with
  \[
    H \vdash\; \thFunTFor(e) = \reify{\code{G}}.
  \]

\item \textbf{diagFTarget} (precomputed):
  $H \vdash\; \cstf(\hat{N}) = \reify{\code{G}}$.

\item \textbf{Chain} steps 2 and 3 through
  $\reify{\code{G}}$ as middle term (trans + sym):
  \[
    H \vdash\; \thFunTFor(e) = \cstf(\hat{N}).
  \]
  The term $\reify{\code{G}}$ appears only as the
  middle of a transitivity chain --- it is never touched by inst.

\item \textbf{Instantiate} $G$ with $e$ (inst on~$x_0$):
  \[
    H \vdash\;
    \TreeEq(\thFunTFor(e),\;
    \mathrm{substTFor}(\hat{N})) = \falseT.
  \]

\item \textbf{Replace} $\thFunTFor(e)$ with
  $\cstf(\hat{N})$ using $\mathrm{cong}_L$ and the chain
  from step~4:
  \[
    H \vdash\;
    \TreeEq(\cstf(\hat{N}),\;
    \mathrm{substTFor}(\hat{N})) = \falseT.
  \]

\item \textbf{Instantiate} $x_{11}$ and $x_{12}$ (inst, using
  the cached computation from Section~4):
  the open $\mathrm{substTFor}(\hat{N})$ becomes
  $\cstf(\hat{N})$:
  \[
    H \vdash\;
    \TreeEq(\cstf(\hat{N}),\;
    \cstf(\hat{N})) = \falseT.
  \]

\item \textbf{Reflexivity} ($\TreeEq(t,t) = \trueT$ via \SchemaF):
  \[
    H \vdash\;
    \TreeEq(\cstf(\hat{N}),\;
    \cstf(\hat{N})) = \trueT.
  \]

\item \textbf{Contradiction}: by trans,
  $H \vdash\; \trueT = \falseT$, violating consistency.
\end{enumerate}

%% ===============================================================
\section{Proof of G\"odel II}

\subsection*{The internal consistency statement}

The code of the contradiction $\trueT = \falseT$ is
\[
  \code{\bot} \;=\; \code{\trueT = \falseT}
  \;=\; \pair{\code{O}}{\code{\pair{O}{O}}}.
\]

\begin{definition}
The \emph{internal consistency statement} is the equation
\[
  \mathrm{Con} \;\equiv\quad
  \TreeEq\bigl(\thFunTFor(x),\;\reify{\code{\bot}}\bigr) \;=\; \falseT
\]
where $x$ is a free variable, implicitly universally quantified.
Since $\TreeEq(a, b) = \falseT$ means $a \neq b$,
$\mathrm{Con}$ says:
\emph{for every proof tree~$x$, the equation encoded by~$x$
is not $\trueT = \falseT$} --- i.e., the system does not
prove $\trueT = \falseT$.
\end{definition}

\noindent
When the hypothesis is a trivially true equation such as
$\trueT = \trueT$, this corresponds exactly to Guard's
Theorem~14: ``$\mathrm{th}(y) \neq \text{``}0 = 1\text{''}$
is valid but unprovable in BRA''~\cite{guard1963}.

\subsection*{Statement}

\begin{theorem}[G\"odel~II, constructive form --- \texttt{conToBot} in \texttt{GodelIIFull.agda}]
\label{thm:godel2cf}
For any hypothesis equation~$H$:
\[
  H \vdash \mathrm{Con} \;\to\;
  H \vdash \trueT = \falseT.
\]
\end{theorem}

\noindent
That is: if the system derives $\mathrm{Con}$,
then it derives $\trueT = \falseT$.  This is the direct constructive
content --- no consistency assumption, no Markov's principle.

\begin{corollary}[G\"odel~II, standard form --- \texttt{godelII} in \texttt{GodelIIFull.agda}]
\label{thm:godel2s}
For any hypothesis equation~$H$:
\[
  \mathit{Consistent}(H) \;\to\;
  H \vdash \mathrm{Con} \;\to\; \bot.
\]
\end{corollary}

\begin{remark}
The constructive form (Theorem~\ref{thm:godel2cf}) is stronger
than the standard form (Corollary~\ref{thm:godel2s}).
Classically they are equivalent by contraposition.
Constructively, our proof directly produces the derivation
$H \vdash \trueT = \falseT$ from $H \vdash \mathrm{Con}$
without any classical reasoning: the consistency
assumption is applied only at the final step to obtain~$\bot$.
\end{remark}

\subsection*{Proof}

The proof has three components.

\medskip\noindent
\textbf{Step 1: Pre-compute the encoding of the G\"odel~I argument.}

\noindent
Specialize the G\"odel~I proof to the hypothesis
$H_0 = G$ (the G\"odel sentence itself).
With this hypothesis, $\mathrm{ruleHyp}$ gives
$G \vdash G$, and the G\"odel~I argument
(Theorem~\ref{thm:godel1cf}) produces a derivation:
\[
  d_\bot : G \vdash \trueT = \falseT.
\]
Apply Theorem~14 (\texttt{thm14E}) to~$d_\bot$ to obtain
the proof encoding $\mathit{enc}_\bot$ and the correctness proof:
\[
  \mathit{corrBot} : \thFunTFor(\mathit{enc}_\bot) = \reify{\code{\bot}}.
\]
Crucially, $\mathit{corrBot}$ is \emph{polymorphic in the hypothesis}:
it is a derivation using only computation axioms and structural
rules, not the hypothesis rule.  Therefore it holds in any
hypothesis context, not just under~$G$.

The term $\mathit{enc}_\bot$ is a specific closed term computed
at the Agda meta-level.  It encodes the entire G\"odel~I
proof: the diagonal (\texttt{diagFTarget}), tree equality
self-comparison (\texttt{treeEqSelf}), variable instantiation,
and all 26~case verifications through \texttt{thm14E}.

\medskip\noindent
\textbf{Step 2: Instantiate the consistency hypothesis.}

\noindent
Given $H \vdash \mathrm{Con}$, apply variable
substitution (\texttt{ruleInst}) with $x := \mathit{enc}_\bot$:
\[
  H \vdash \TreeEq(\thFunTFor(\mathit{enc}_\bot),\;\reify{\code{\bot}}) = \falseT.
\]

\medskip\noindent
\textbf{Step 3: Derive the contradiction.}

\noindent
From $\mathit{corrBot}$ (polymorphic, instantiated at~$H$):
\[
  H \vdash \thFunTFor(\mathit{enc}_\bot) = \reify{\code{\bot}}.
\]
Replace $\thFunTFor(\mathit{enc}_\bot)$ with $\reify{\code{\bot}}$
by congruence:
\[
  H \vdash \TreeEq(\reify{\code{\bot}},\;\reify{\code{\bot}}) = \falseT.
\]
But $\TreeEq(t, t) = \trueT$ for all~$t$ (by \SchemaF{}):
\[
  H \vdash \TreeEq(\reify{\code{\bot}},\;\reify{\code{\bot}}) = \trueT.
\]
By transitivity: $H \vdash \trueT = \falseT$,
contradicting $\mathit{Consistent}(H)$.
\qed

%% ===============================================================
\section{Discussion}

\subsection*{Comparison with Guard and Rose}

Guard~\cite{guard1963} proves G\"odel~II (his Theorem~14) by
explicitly constructing a functor~$h(x)$ within BRA such that
$\mathrm{th}(h(x)) = \text{``}0 = 1\text{''}$ whenever
$\mathrm{th}(x) = j$ (the G\"odel sentence code).  This uses
his Theorem~13 (internalization: $f(x) = y \supset
\mathrm{th}(D_f(x)) = \text{``}fx = y\text{''}$).

Rose~\cite{rose1961} proves G\"odel~II (his Theorem~18) by
applying his Theorem~16 (internalized Schema~F) to the function
$a(\mathrm{se}(\mathrm{nu}(N), N), \mathrm{th}(y))$, using the
consistency hypothesis $\mathrm{th}(x) \neq e(\mathrm{th}(z))$.
Cook~\cite{cook1975} points out a possible error in Rose's
Theorem~16 (p.~134 of~\cite{rose1961}): the function $f(s(x))$
is neither identically zero nor identically non-zero, which
invalidates the argument at that step.

Our proof takes a different route: rather than constructing~$h$
as an object-level functor or applying internalized Schema~F,
we \emph{pre-compute} the encoding $\mathit{enc}_\bot$ at the
Agda meta-level from the concrete G\"odel~I derivation.
The polymorphicity of Theorem~14's correctness proof
($\mathit{corrBot}$ holds in any hypothesis context) is the
key property that makes this work, bypassing Rose's
problematic Theorem~16 entirely.

\subsection*{The role of Theorem~14}

The entire non-trivial content of the proof resides in
Theorem~14 (\texttt{thm14E}).  The 26~case files verify
that $\thFunTFor$ correctly tracks each derivation rule.
The G\"odel~II proof applies Theorem~14 to the specific
derivation~$d_\bot$ (the G\"odel~I argument with hypothesis~$G$),
obtaining an encoding whose $\thFunTFor$ image is the
contradiction code.  The instantiation of $\mathrm{Con}$
at this encoding then produces the contradiction.

\subsection*{Connection with Cook's PV}

Cook~\cite{cook1975} suggests (Section~8) that a natural
formalization of feasibly constructive proofs should take a
\emph{programming approach}, where ``proving $f(x) = g(x)$
amounts to proving the equivalence of two programs.''
The Guard/Rose system realizes this vision literally:
proving $f(x_0) = g(x_0)$ via \SchemaF{} \emph{is} proving
that two recursive programs (the functors $f$ and~$g$) compute
the same function.
Every derivation in the system is a chain of equational
rewrites on functor expressions --- no quantifiers, no
propositional connectives, only computation axioms and
the uniqueness-of-recursion principle.

%% ===============================================================
\section{Correspondence with Rose}

\begin{center}
\begin{tabular}{ll}
\hline
\textbf{Rose~\cite{rose1961}} & \textbf{This formalization} \\
\hline
$th(y)$ & $\thFunTFor$ \\
$se(x,y,z)$ / $st(x,y,z)$ & $\mathrm{substTFor}$ / $\mathrm{closedSubstTFor}$ \\
$nu(x)$ & implicit via $\hat{N}$ \\
$N$ (template code) & $\mathrm{templateCode}$ \\
Schema~F & $\mathrm{ruleF}$ (uniqueness of recursion) \\
Theorem~10 & $\mathrm{thm14E}$ ($\mathrm{ProofE}$) \\
Theorem~14 & $\mathrm{thm14E}$ \\
Theorems~15--16 & Bypassed (see Discussion); cf.~\cite{cook1975} on Thm~16 \\
Theorem~17 & $\mathrm{diagFTarget}$ \\
Theorem~18 & \texttt{godelII} in \texttt{GodelII.agda} (= G\"odel I) \\
--- & \texttt{godelII} in \texttt{GodelIIFull.agda} (= G\"odel II) \\
\hline
\end{tabular}
\end{center}

%% ===============================================================
\section{File structure}

\begin{center}
\begin{tabular}{lrl}
\hline
\textbf{File} & \textbf{Time} & \textbf{Role} \\
\hline
Base.agda & $<1$s & Natural numbers, equality, sigma types \\
Term.agda & $<1$s & Terms, functors, coding functions \\
Step.agda & $<1$s & Derivation system (26 rules + \SchemaF) \\
StepReduce.agda & $<1$s & Derived reduction lemmas \\
ThFun.agda & $<1$s & Metalevel proof evaluator \\
ThFunCorrect.agda & $<1$s & Metalevel correctness \\
ThFunTFor.agda & $<1$s & Object-level proof evaluator (26 branches) \\
ThFunTForDefs.agda & $<1$s & Definitions for thFunTFor cases \\
ThFunTForCases0--3.agda & $<1$s & Per-tag case proofs \\
ThFunTForCorrectDefs.agda & $<1$s & Correctness infrastructure \\
Thm14E.agda & $<1$s & Theorem~14 (1402 lines, all 26 cases) \\
SubstTFor.agda & $<1$s & Coded substitution functor \\
SubstCorrect.agda & $<1$s & Substitution correctness \\
SubstTForCorrect.agda & $<1$s & $\mathrm{closedSubstTFor}$ correctness \\
SubstTForPrecomp.agda & 108s & \textbf{Heavy}: precomputed abstractions \\
SubstTForPrecomp2.agda & $<1$s & Assembly from opaque pieces \\
NdDispatchProofs.agda & $<1$s & Node-dispatch case proofs \\
PairPassthrough.agda & $<1$s & Pair passthrough lemmas \\
PairPassthroughAll.agda & $<1$s & Extended passthrough \\
IntermediatePassthrough.agda & $<1$s & Intermediate passthrough \\
ExtractorRed.agda & $<1$s & Extractor reduction lemmas \\
GodelII.agda & $<1$s & G\"odel I (168 lines) + treeEqSelf \\
ProofEAny.agda & $<1$s & mkProofEAny + strengthened G\"odel I \\
GodelIIFull.agda & $<1$s & \textbf{G\"odel II} ($\sim$90 lines) \\
CookPV.agda & $<1$s & Boolean ops, De Morgan, excluded middle \\
NoPChecker.agda & $<1$s & No proof-checker (corollary of G\"odel II) \\
\hline
\end{tabular}
\end{center}

All results compile with Agda~2.9.0 using
\texttt{-{}-safe -{}-without-K -{}-exact-split}.
No postulates.  The 108-second computation is cached and not repeated.

\begin{thebibliography}{9}
\bibitem{guard1963}
  J.~R.~Guard,
  \emph{Lecture Notes on Recursive Arithmetic},
  Princeton University, 1963.

\bibitem{rose1961}
  H.~E.~Rose,
  ``On the consistency and undecidability of recursive arithmetic,''
  \emph{Zeitschr.\ f.\ math.\ Logik und Grundlagen d.\ Math.},
  vol.~7, pp.~124--135, 1961.

\bibitem{cook1975}
  S.~A.~Cook,
  ``Feasibly constructive proofs and the propositional calculus,''
  \emph{Proceedings of the Seventh Annual ACM Symposium on Theory of Computing},
  pp.~83--97, 1975.
\end{thebibliography}

\end{document}
